\documentclass[12pt,a4paper,oneside, reqno]{amsart}
%\documentclass[12pt]{book}%{amsart}
\usepackage[top=.65in, bottom=0.85in, left=0.8in, right=1.2in]{geometry}
\usepackage[hyphens]{url}
\usepackage{graphicx}
\usepackage{fancyhdr}
%\usepackage{fullpage}
\usepackage{listings}
\usepackage{multirow}
\usepackage{color,soul}
\usepackage{changes}
\usepackage{lipsum}
\usepackage{enumitem}
\usepackage{ulem}
\usepackage{soul}

\usepackage[backend=biber,style=apa]{biblatex}
\addbibresource{bibliography.bib}

%%%%%%%%%%%%%%%%%%%%%%%%%%%%%%%%%%%%%
\newtheorem{theorem}{Theorem}[section]
\newtheorem{lemma}[theorem]{Lemma}
%\theoremstyle{definition}
\newtheorem{definition}[theorem]{Definition}
\newtheorem{example}[theorem]{Example}
\newtheorem{xca}[theorem]{Exercise}

\setlength\parindent{0pt}


\fancyfoot[R]{\thepage}

%\addtolength{\oddsidemargin}{-0.20in}
%\addtolength{\evensidemargin}{-0.20in}
\addtolength{\textwidth}{0.6in}
%\addtolength{\topmargin}{1.0in}
%\addtolength{\textheight}{1.0in}

\setlength{\parskip}{4pt}

\begin{document}

\newpage

\thispagestyle{empty}


\mbox{}

\vspace{1pt}

\begin{flushright}


{\large Response to Reviewers}

\end{flushright}

\vspace{15pt}


\setlength{\leftskip}{0cm} We would like to express our gratitude for your helpful comments. We appreciate the recognition of the novelty of our study, scale of the dataset and importance of our findings for understanding post-pandemic urban resilience and socioeconomic equity in Latin America. We have made changes to the manuscript in response to the comments. As a result, we have enhanced the robustness of the analysis, revised the figures, and updated the relevant sections of the manuscript for increased clarity. 

\vspace{10pt}

\textbf{REVIEWER 1}
\vspace{-16pt}
\begin{center}
    \rule{\textwidth}{.1pt}
\end{center}

\vspace{5pt}

\setlength{\leftskip}{0cm} \textit{\textbf{Reviewer 1 (R1)}: The present work addresses an important question. The findings are novel, the dataset is large, and the implications for urban resilience and equity are important. However, the robustness of the results is currently undermined by unresolved concerns about data reliability, adjustment validation, and confounding. Without stronger evidence that the patterns are not artifacts of Meta-Facebook’s processing or of model assumptions, the paper risks overstating its conclusions. For these reasons, I recommend the manuscript to go through a major revision cycle.}

\setlength{\leftskip}{1cm}\textbf{Authors (A):} We appreciate your comments and concerns. We have taken them all onboard to improve the manuscript and analysis. Please see below our responses describing the amendments and additional data analysis that we have undertaken.

\setlength{\leftskip}{0cm}\textit{\textbf{R1, Comment 1 (R1.1)}: The reliability of Meta-Facebook mobility data is questionable. The dataset is large-scale but raises substantial concerns. As mentioned in the Method Section, Meta-Facebook applies multiple privacy-preserving techniques, including random noise injection, spatial smoothing, and suppression of data count below 10. These interventions may distort observed patterns, particularly in low-density areas. The claim that low-density and deprived areas quickly returned to baseline may be an artifact of these processing rules.}

\setlength{\leftskip}{1cm}\textbf{A:} We share the reviewer’s concerns about the potential influence of Meta’s privacy-preserving adjustments on data reliability in low-density areas. At the same time, we acknowledge that these adjustments are necessary for ethical and privacy reasons. Because of the same concerns raised by the reviewer, we did undertake measures to mitigate potential distortions resulting from the described adjustments through careful data imputation and validation. We detailed all the steps we took in the Methods and Supplementary Information. We revised these sections to address explicitly the concerns raised by the reviewer on small numbers and low population-density areas. Here, we provide a brief summary of the most crucial steps:
\begin{itemize}[noitemsep, topsep=0pt]
    \item Validation. We assessed the degree of correspondence between the spatial distribution of Facebook active user counts during the baseline period against WorldPop population estimates. The two datasets show a statistically significant and strong positive correlation ($R = 0.88–0.98$, $p<0.05$) for the three countries in our sample - see Supplementary Figures (SF) 2-4. These results indicate that Facebook data adequately capture the spatial patterns of population density across sparsely populated areas and densely urban areas (i.e. along the rural-to-urban hierarchy) in the three countries. 
    \item Imputation. Based on the validated baseline population data, we calibrated a spatial interaction model that predicts baseline mobility flows as a function of population at origin and destination locations and their distance. This model was then used to estimate missing flows suppressed by Meta’s privacy thresholds (i.e. records with counts under 10) during the ``crisis'' period. In Supplementary Figure SF5, we show that our model predicted flows and the observed baseline flows have a remarkable patterns of correspondence, validating our imputation approach.
\end{itemize}

\setlength{\leftskip}{1cm}\textcolor{blue}{\textbf{Action:} We revised the text in the ``Results'' section to more clearly explain how the adjustment procedures that Meta implements may distort their data and the approaches we undertook to mitigate these issues and validate the results. Specifically, we added the following highlighted text: \textit{``\uline{We know that Facebook suppresses population counts below 10 to protect users' privacy and that user engagement declined over time in our dataset. These adjustments may influence the observed patterns of mobility, especially in sparsely populated areas. To mitigate these issues, we applied a series of data corrections (see Figure 1). First, we modelled the Facebook user population data as a function of Worldpop population data to
generate a complete spatial population surface. Second, we modelled the Facebook mobility data to estimate origin-destination
flows that were suppressed by the Facebook algorithm. Third, we used the temporal median of the sum of active user counts
divided by the sum of active user counts across all spatial units on a given day to normalise for fluctuations in Facebook user engagement over time. Full details of these procedures are provided in the Methods section, and ST1 summarises the extent of
data imputation.}''}}

\setlength{\leftskip}{0cm} \textit{\textbf{R1.2}: Coverage bias is another concern. Meta-Facebook users may not be evenly distributed across regions or socio-economic groups. To strengthen confidence in the findings, the paper should incorporate ground truth data sources. Even a partial validation exercise would add credibility.}

\setlength{\leftskip}{1cm}\textbf{A:} We agree that coverage bias represents an important concern of Meta–Facebook data, as users may not be distributed across regions or socio-economic groups in a way that reflects the underlying populations. At the same time, we acknowledge that, for privacy reasons, Meta data do not include any individual-level sociodemographic or socioeconomic attributes, which prevents the application of standard post-stratification or re-weighting techniques. Fully correcting for representation bias in this dataset is technically complex. It first requires establishing a comparable benchmark population and mobility flows or ground truth dataset; that is, capturing actual population counts every eight hours, every day from March 2020 to May 2025. Such population data do not exist, but this is a key input requirement to measure the extent of coverage bias, to then measure it and seek for appropriate algorithm to correct it. We believe that a carefully designed methodology to appropriately adjust for bias in mobile-phone-derived population and human mobility data is yet to be developed. In fact, we have a funded project which aims to develop such methodology. We recognise that census population data are normally used to correct for coverage bias. However, this has been shown to regress mobile-phone-derived population counts back to be representative at the census date. Yet, this defeats the essential point of using mobile-phone-derived data - which is to render a picture of population dynamics for which census data are not available. Hence, direct demographic weighting is not feasible in our context. 

Instead, as suggested by the reviewer, we undertook a partial validation exercise. We assessed the representativeness of the Facebook data by comparing the percentage distribution of population counts across population density and Relative Deprivation Index (RDI) categories with the corresponding distributions derived from WorldPop data. We treat WorldPop dara as ground truth. This validation assessment is shown in the new Supplementary Fig. 8 (SF8). The comparison allows us to highlight any potential under- or over-representation of specific population groups in Facebook data. Our findings indicate that the distributions are closely aligned with Worldpop across all three countries (especially after data imputation) suggesting that Facebook data reasonably approximates population structures along rural–urban and socio-economic dimensions.

\textcolor{blue}{\textbf{Action:} In particular, we have taken three actions to address concerns on coverage bias. 
\begin{itemize}[noitemsep, topsep=0pt]
    \item First, we have produced a new Supplementary Figure (SF8) showing the population share in each population density and Relative Deprivation Index category by country, based on Meta–Facebook data and WorldPop data, which we treat as the ground truth. The results show that the Facebook data are reasonably representative, especially after imputation, as their distributions across population density and Relative Deprivation Index categories closely align with those derived from WorldPop.
    \item Second, we revised and expanded the ``Methods'' section, specifically the ``Classification of tiles by urbanisation and socio-economic deprivation levels'' subsection. We do this to describe the results of our validation exercise described above which is reported in the new Supplementary Figures SF8. 
    \item Third, we have revised our models and incorporated new population categories into the modelling of time trends, to account for (i) coverage bias, as they capture demographic patterns that may influence Facebook user representation (see \parencite{cabrera2025bias}), and (ii) differences in mobility behaviours and capacity for mobility recovery. Particularly we included the percentage of elderly population within each spatial unit using Worldpop data. We provide further details about the inclusion of these data in relation to comments R1.4 and R2.8.
\end{itemize}}

\setlength{\leftskip}{0cm} \textit{\textbf{R1.3}: The authors apply several non-trivial data adjustments (imputation of missing counts, calibration with WorldPop, and correction for declining Facebook engagement). They also use STL decomposition and mixed-effects modeling to estimate recovery trends. However, the robustness of these steps is insufficiently demonstrated. The Methods section provides detailed procedures, but the results would benefit from explicit validation checks that show how findings hold under alternative assumptions.}

\setlength{\leftskip}{1cm}\textbf{A:} We have undertaken additional analysis to offer further evidence the robustness of our approach, demonstrating that our findings hold under various model specifications and method selection.

\textcolor{blue}{
\textbf{Action:} Specifically, we implemented the following changes: 
\begin{itemize}[noitemsep, topsep=0pt]
    \item First, we applied two alternative trend extraction techniques, to evaluate the robustness of the recovery trends: (i) STL decomposition (Seasonal-Trend decomposition using Loess) \parencite{cleveland1990stl} and (ii) a machine learning model developed by Meta, known as Prophet \parencite{prophet17}. The latter estimates adaptive piecewise-linear trends with seasonality. We compared alternative trends from each of these methods to our original trend estimates and found strong correlations ($R = 0.77–1.00$), suggesting that the trends extracted through different methods are consistent and capture similar underlying mobility patterns. We report these results in the new Supplementary Table ST2.
    \item Second, we have implemented an alternative model specification (Model 8) based on a Generalised Additive Model (GAM) framework, to test the robustness of our mixed-effects linear model specifications. The full set of GAM estimates is provided in the revised Supplementary Tables ST5–ST21. The GAM estimates yield lower Akaike Information Criterion (AIC) values than our mixed-effects specifications but have clear signs of overfitting, displaying high effective degrees of freedom associated with the category-specific spline terms. In addition, we find that the mixed-effects models offer a more transparent and interpretable structure for our research questions. We also derived recovery time estimates using the GAM specification, and these are consistent with those obtained from the mixed-effects models, as shown in the revised Supplementary Tables ST22–ST25. 
    \item Third, we expanded the text in the ``Results'' and ``Methods'' sections to discuss the results from the additional analysis.
\end{itemize}
}

\setlength{\leftskip}{0cm} \textit{\textbf{R1.4}: The analysis focuses primarily on population density and socio-economic deprivation. While these are important, other structural covariates could confound the observed differences. Factors such as age distribution, education levels and transport infrastructure availability likely influence both mobility patterns and recovery capacity. Without controlling for these, it is difficult to rule out alternative explanations for the observed gradients. The discussion should also more directly acknowledge these limitations.
}

\setlength{\leftskip}{1cm}\textbf{A:} We agree that multiple covariates may influence both mobility behaviour and recovery dynamics. As such, we analysed variations in these components across deprivation and population density categories. Particularly, we used the Relative Deprivation Index (RDI) because it is a composite index that captures six overlapping dimensions reflecting geographic differences in: (i) built-up density; (ii) child dependency ratio; (iii) infant mortality rate; (iv) human development; (v) nighttime lights intensity; and (vi) nighttime lights slope \parencite{RDI2022}. These dimensions are intended to reflect differences in education, health and infrastructural development \parencite{RDI2022}, mitigating the concern that the analysis omits key structural determinants. 

We also note that (i) our models incorporate a random intercept to account for unobserved structural differences across regions that are not explicitly captured by the included covariates \parencite{gelman2007data}, (ii) our choice of covariates for the models is based on conceptual relevance and on availability of data across countries and at a spatial resolution fine enough to match the granularity of our mobility analysis. 

Within these constraints, we recognise the value of explicitly analysing mobility for areas with varying age structure, as age can plausibly influence mobility patterns \parencite{schwartz1976migration} and capacity to recover, as well as the levels of digital data coverage (as we see in \parencite{cabrera2025bias}). Accordingly, we extended our analysis to control for age groups, as detailed below.

\textcolor{blue}{
\textbf{Action:} We have revised the manuscript taking the following actions:
\begin{itemize}[noitemsep, topsep=0pt]
    \item First, we expanded the modelling framework to incorporate demographic age structure by computing the old-age dependency ratio for each spatial unit, defined as the percentage of residents aged 65 and above relative to the total population. We focus specifically on the old-age dependency ratio, as child dependency is already indirectly captured through the RDI. We derived the old-age dependency ratio from WorldPop data, which provide globally consistent, fine-resolution population estimates to match the spatial resolution of our mobility data. To examine whether demographic composition influences mobility patterns, we categorised spatial units into three groups (low, medium, high) based on their old-age dependency levels. The visual results are collected in the new Supplementary Figure SF11. The statistical models (SF11 and new Supplementary Tables ST4-9) show no significant differences in recovery rates (time slopes) across old-age dependency ratio categories. Intercepts do differ across the three old-age dependency categories, indicating variation in the initial magnitude of the mobility shock, but these differences do not follow a systematic gradient across countries. This contrasts with the reported patterns for population density and RDI categories. For this reason, we focus on population density and RDI for the main analysis. We report the age-structure results in the Supplementary Information.
    \item Second, we have revised the text in the ``Results'' section to (i) describe in greater detail the multidimensional components of the RDI, which capture educational, health and infrastructure characteristics relevant to mobility (further details of the changes in the text can be found in comment R3.7), and (ii) to describe the inclusion of the additional old-age dependency ratio variable in our models.
\end{itemize}
}

\setlength{\leftskip}{0cm} \textit{\textbf{R1.5}: The Discussion section emphasizes mobility disparities and risks of urban center decline, but the societal implications could be more thoroughly explored. For instance: What do persistent declines in urban core mobility mean for equity in access to jobs, education, or healthcare? Are there policy recommendations specific to Latin American governance structures, beyond general calls for inclusive planning? Expanding this discussion would help position the study as not only descriptive but also impactful for urban policy debates.
}

\setlength{\leftskip}{1cm}\textbf{A:} Thanks for pointing out that our work may have wider implications and we should elaborate on them. We have expanded the text to discuss the potential broader societal implications of persistent mobility disparities and urban decline in relation to equitable access to essential services and policy design in Latin America countries. 

\textcolor{blue}{
\textbf{Action:} We have modified and expanded the ``Discussion''. The final two paragraphs now explicitly focus on the implications of our findings for (i) social equity, highlighting how uneven mobility recovery may reinforce existing socio-economic inequalities and spatial mismatches between residence and access opportunity, and (ii) urban policy, with a stronger emphasis on the specific governance and fiscal constraints of Latin American metropolitan regions. We also discuss targeted, context-specific policy responses and the role of near real-time mobility data in supporting adaptive, evidence-based urban planning. We copied the new text in the response to comment R3.2.
}

\setlength{\leftskip}{0cm} \textit{\textbf{R1.6}: The Figure presentation could be improved. For example, Figure 1 outlines data inputs and adjustment techniques. However, it is dense and difficult to follow. Figure 2’s boxplots are not fully intuitive. It is difficult to see at a glance which groups experienced larger sustained declines.
}

\setlength{\leftskip}{1cm}\textbf{A:} We have revised Figures 1 and 2 to make their key messages more intuitive.

\textcolor{blue}{
\textbf{Action:} We have revised Figure 1 and 2. For Figure 1, we have simplified the visualisation, keeping only the essential elements; that is, the input data, the data calibration process, the analysis and modelling and the validation of outputs. We have emphasised more the output validation section. We have updated the caption to better describe the overview of our methodology, and direct the readers to the ``Methods'' section for more detailed descriptions. For Figure 2, we have separated the boxplots from the year 2020 and 2022 into two separate rows. We have also assigned different panels for analysis by population density categories and by relative deprivation index. We have updated the colour palette for the relative deprivation categories, to better distinguish it from the population density categories. We have also modified the other Figures according to the suggestions given by reviewer 2 (please see comments R2.9 and R2.10).
}

\vspace{10pt}

\setlength{\leftskip}{0cm} \textbf{REVIEWER 2}
\vspace{-16pt}
\begin{center}
    \rule{\textwidth}{.1pt}
\end{center}

\vspace{5pt}

\setlength{\leftskip}{0cm} \textit{\textbf{Reviewer 2, Comment 1 (R2.1)}: While the methodological approach is well thought out, I regret to say that I cannot recommend the paper for publication in its current form due to a number of significant limitations in the data and methods. Below, I detail my concerns and offer suggestions for improvement:
}

\setlength{\leftskip}{1cm}\textbf{Authors (A):} We appreciate the concerns raised about data limitations and methodological robustness. Many of the issues raised align with those highlighted by Reviewer 1. We have undertaken a series of substantial revisions to address them, as we outline in this document.

\setlength{\leftskip}{0cm} \textit{\textbf{R2.2}: In the title, the study claims to focus on Latin America but includes only three countries (Argentina, Colombia, and Chile), omitting key countries like Brazil and Mexico.}

\setlength{\leftskip}{1cm}\textbf{A:}  We acknowledge that our analysis focuses on three Latin American countries and does not include the entire region. Generally, paper titles do not include the geographical area from where they source the data, even if the focus is on one city or country. Our use of the phrase ``in Latin America'' in the title is explicit and reflects the international context and relevance of the findings. The specific countries studied are clearly identified in the Abstract. The selected countries were chosen to represent diverse economic, institutional and spatial contexts within Latin America. If possible, we would prefer to retain the current title, as we believe it accurately captures the conceptual and geographical focus of the paper. It also keeps it concise and accessible to a general readership.

\setlength{\leftskip}{0cm} \textit{\textbf{R2.3}: Before addressing the first research question, the study should first verify whether the COVID-19 pandemic led to observable changes in urban mobility patterns as captured by cell phone data in three Latin American countries.}

\setlength{\leftskip}{1cm}\textbf{A:}  We agree that establishing whether the COVID-19 pandemic led to observable changes in mobility is a necessary first step. Prior work has already established such changes using the same (Meta–Facebook) and alternative data sources (i.e. Google mobility data), including several of our studies (e.g. \parencite{rowe2023a, covidMexico24, demography25}). There is consistent evidence for sharp reductions in mobility during the onset of the pandemic across Latin American cities. In the present paper, our focus is to analyse the persistence and inequality of post-pandemic recovery of mobility patterns across socio-economic and population density gradients, rather than focusing on the mobility declines induced by the pandemic in the first half of 2020.

\textcolor{blue}{
\textbf{Action:} We added a statement in the ``Introduction'' to establish this context i.e. that mobility declined during the early stages of the COVID-19 pandemic. The following text was added: \textit{\uline{``These data sources have consistently documented sharp and widespread mobility declines during the early months of the pandemic across Latin American cities} \parencite{RoweCCA24, covidMexico24, elejalde2024, demography25}\uline{, providing a well established empirical foundation for focusing on the persistence and inequality of the post pandemic recovery in the present study.''}}.
}

\setlength{\leftskip}{0cm} \textit{\textbf{R2.4}: The study’s primary limitation is in the data and temporal framing. The baseline period used for comparison consists of only 45 days (Jan 24 – Mar 10, 2020), which does not adequately capture typical mobility patterns or account for pre-pandemic variability. And this reference period is unbalanced in terms of day-of-week representation (some days of week have 7 samples while others only have 6, which will affect the standard deviation of the benchmark), potentially introducing systematic bias, especially for weekly cyclical travel patterns. Data collection begins around the onset of the pandemic, without a sufficiently long pre-COVID baseline, which limits the study’s ability to accurately assess the extent of disruption and recovery trajectories. Besides, the coarse temporal resolution (8-hour intervals) limits the validity of any inference about daily commuting patterns. For example, in Argentina and Chile, time windows such as 5:00-13:00 or 13:00-21:00 cannot simultaneously capture both morning and evening peak periods, making the identification of ``daily commutes'' problematic.}
 
\setlength{\leftskip}{1cm}\textbf{A:} Thanks for the thoughtful comment. Meta provides a baseline period extending for 45 days before March 2020. Hence the existing data available from Meta prevent us from having a longer pre-pandemic baseline. However, first, we would like to argue that a baseline period of 45 days is completely reasonable. For Malaria, for example, the official operational guide for disease surveillance published by the World Health Organization (WHO) establishes a baseline period of three weeks (21 days) \parencite{WHO2025:UNKNOWNKEY} prior to the onset of the disease. Also, our baseline period extends from Jan 24 to Mar 10, 2020. This period did not overlap with mobility restrictions in the countries in our sample. Restrictions started on March 19th in Argentina, 25th in Colombia and 26th in Chile. The Meta mobility data thus effectively capture mobility patterns pre-pandemic.

Second, we would also like to recognise the exceptionally high temporal resolution of the Facebook mobility data relative to conventional administrative data sources, which are typically available only through travel surveys or censuses conducted every several years, as well as the density of observations in each time window. Other mobile phone app sources only capture a few thousands of devices each day in countries like the UK. Additionally, while the data may not capture daily commuting peaks, they capture general human mobility for different purposes representing a measure of vitality of communities which is what we intend to capture. We thus believe that the daily observations provided by the Meta dataset allow for a much more granular and dynamic assessment of mobility evolution over time than would otherwise be possible based on existing data sources available in the Latin American countries in our analysis.

\textcolor{blue}{
\textbf{Action:} We have taken three actions.
\begin{itemize}[noitemsep, topsep=0pt]
    \item First, we have revised the text to clarify that the baseline data is provided by Meta-Facebook. Particularly, we have added in the ``Methods - Data - Facebook Movement During Crisis'' section the following: \textit{``In addition, each dataset includes baseline movement counts, which represent the estimated number of people moving before COVID-19. These baseline counts \ul{provided by Meta} are calculated based on a 45-day period ending on March 10, 2020, using the average for the same time of day and day of the week within this period. For example, Meta calculates the baseline for the 00:00–08:00 (PT) window on Mondays by averaging all Monday records for that time window across the 45-day period. Further details about the baseline can be found in} \parencite{Maas19}.''}. We have also amended the ``Methods - Data - Facebook Movement During Crisis'' section to clarify that \textit{``Similar to the Facebook Movements datasets, the Populations datasets include baseline counts, calculated by averaging values for the same day of the week and time of day over the 45-day pre-pandemic \ul{baseline period provided by Meta.}''}
    \item Second, we have clarified that our focus was the analysis of ``daily mobility'' as a proxy of community vibrancy or vitality, rather than ``commuting'' patterns, as we agree that the chosen time window cannot effectively capture morning and/or evening peak periods. Specifically, we have updated the ``Meta-Facebook data pre-processing'' section within ``Methods - Data'', where we have changed \textit{``We only included Facebook data for one time window per day: 08:00–16:00 Pacific Time (PT) corresponding to 13:00–21:00 in Argentina and Chile, and 11:00–19:00 in Colombia. This time window is used to capture daily commutes.''} to \textit{``We only included Facebook data for one time window per day: 08:00–16:00 Pacific Time (PT) corresponding to 13:00–21:00 in Argentina and Chile, and 11:00–19:00 in Colombia. \ul{This time window is used to capture daily mobility.}''}
    \item Third, we have added a paragraph in the ``Discussion'' describing the limitations of our study (see response to comment R3.11 for full details), where we explicitly acknowledge the constraints of the pre-pandemic baseline period provided by Meta and discuss its implications for estimating absolute mobility changes.
\end{itemize}
}


\setlength{\leftskip}{0cm} \textit{\textbf{R2.5}: Depending on extensive data imputation, especially in suburban areas as shown in Figure 1, I have serious concerns about the integrity of the derived mobility metrics. While missing data are inevitable in mobility datasets, the degree of imputation in this study appears substantial. Yet, the authors do not provide sufficient transparency regarding the percentage of data imputed, and nor its spatial distribution. Because much of the study's findings and conclusions are derived from these imputed values, more assessment of their validity is needed and I strongly recommend including summary statistics and robustness checks.}
 
\setlength{\leftskip}{1cm}\textbf{A:} We included detailed information on the extent of imputed data in Supplementary Table 1 and on the spatial distribution of these imputations in Supplementary Figure SF1. We also extended the validation procedures for the data imputation, by assessing the degree of correspondence between the spatial distribution
of Facebook active user counts during the baseline period against WorldPop population estimates, and by confirming that our predicted populations and flows show a high degree of correspondence with the observed ones (see SF2-SF5).

\textcolor{blue}{
\textbf{Action:} Specifically, we have implemented the following changes to the manuscript.
\begin{itemize}[noitemsep, topsep=0pt]
    \item First, we have revised Supplementary Table S1 and associated text in the main manuscript (in the ``Results'' and ``Methods'' sections) and Supplementary Information to describe the extent of imputed data more prominently.
    \item Second, we provide extensive validation procedures for the imputed data, as described in the response to comments R1.1, R1.2 and R2.7. In particular, the newly revised SF8 represents clearly the extent of representativeness of the Facebook data before and after data imputation. Our imputation procedure improves the representativeness of the data across rural-urban and socio-economic categories. Figure SF5 also shows a remarkable degree of correspondence between the imputed flows and the observed flows.
\end{itemize}
}

\setlength{\leftskip}{0cm} \textit{\textbf{R2.6}: The manuscript assumes a quadratic form for the trend component used to model mobility recovery. However, recovery from a systemic shock like COVID-19 often follows non-linear and non-symmetric patterns, which may be better captured by spline models, piecewise regressions, or other flexible approaches. The authors should justify the use of a quadratic model and compare it with alternatives.}
 
\setlength{\leftskip}{1cm}\textbf{A:}  Our use of a quadratic form to model the trend was motivated by parsimony and interpretability. It allows us to capture non-linear recovery trajectories, without imposing excessively complex functional forms and overfitting the data. This approach also allow us to compute recovery times analytically. Nonetheless, we recognise that recovery from a systemic shock, such as COVID-19 may display more complex dynamics. We therefore consider alternative model specifications to accomodate more complex, non-linear functional forms.

\textcolor{blue}{
\textbf{Action:} We assessed the following approaches:
\begin{itemize}[noitemsep, topsep=0pt]
    \item First, we used a Generalised Additive Model (GAM) framework and tested an alternative model specification (Model 8) to assess non-linear recovery dynamics. The GAM specifications achieve lower Akaike Information Criterion (AIC) values, but they also show evidence of overfitting (high effective degrees of freedom). Despite this, we estimated the recovery times (now reported in ST22-ST25) based on these spline-based models and the results are consistent with those from our quadratic mixed-effects models. Overall, we believe that our original quadratic form offers a stable and interpretable approximation of the underlying trend. We therefore report these results in the main text, and the alternative specifications and comparative analysis in the newly revised Supplementary Tables ST5–ST21.
    \item Second, we have revised the ``Results'' and ``Methods'' sections to describe the GAM-based robustness analysis and report the consistency of recovery time estimates across model specifications.
\end{itemize}
}

\setlength{\leftskip}{0cm} \textit{\textbf{R2.7}: While the authors briefly acknowledge limitations of cell phone data, the demographic biases inherent in such data (e.g., underrepresentation of older populations or those with limited smartphone access) are not adequately addressed. More discussion is needed regarding how these biases may systematically affect the estimation of mobility trends in marginalized populations.}
 
\setlength{\leftskip}{1cm}\textbf{A:} We agree that demographic bias represents an important concern of mobility data derived from mobile phone location data, such as those used in our analysis. Users are not always distributed across regions or socio-economic groups in a representative way. However, Meta data does not include any individual-level sociodemographic or socioeconomic attributes of users for privacy reasons, which prevents the application of standard post-stratification or re-weighting techniques. Consequently, fully correcting for representation bias in this dataset is not possible with the existing data. Instead, we assessed the representativeness of Meta-Facebook data. To this end, we compared the geographic distribution of population counts across population density and Relative Deprivation Index (RDI) categories based on Meta-Facebook versus WorldPop data (after applying our imputation corrections). We report this analysis in the new Supplementary Figure SF8. The results provide evidence of a high degree of correspondence between the Meta-Facebook users data and WorldPop population data across areas for all three countries. Overall, these results indicate that the Meta-Facebook data offer a reasonable approximation of population structures along the rural–urban and socio-economic geography of the countries in our sample.

\textcolor{blue}{
\textbf{Action:} We have taken the following three actions to address concerns on coverage bias:
\begin{itemize}[noitemsep, topsep=0pt]
    \item First, we have produced a new Supplementary Figure (SF8) showing the population share in each population density and relative deprivation index category by country, based on (i) Meta-Facebook data and (ii) Worldpop data as an external dataset. The results demonstrate that Facebook data offer a relatively good level of representativeness by comparing their distribution across population density and relative deprivation index categories to similar Worldpop data.
    \item Second, we revised and expanded ``Methods'' section, specifically the ``Classification of tiles by urbanisation and socio-economic deprivation levels'' subsection. We do this to describe the results of our validation exercise described above which is reported in the new Supplementary Figures SF7 and SF8.
    \item Third, we have revised our models and incorporated new population categories into the modelling of time trends, to account for (i) coverage bias, as they capture demographic patterns that may influence Facebook user representation (see \parencite{cabrera2025bias}), and (ii) differences in mobility behaviours and capacity for mobility recovery. Particularly we included the percentage of elderly population within each spatial unit using Worldpop data. We provide further details about the inclusion of these data in relation to comments R1.4 and R2.8.
\end{itemize}
}

\setlength{\leftskip}{0cm} \textit{\textbf{R2.8}: The analysis includes a deprivation metric but omits other potentially informative socioeconomic indicators (e.g., income, employment, education) that could be derived from national census data (e.g. 2022 Argentina census data), even though in some countries these data may be more than five years out of date. Incorporating these variables can enhance the study's depth and help disentangle the multifaceted drivers of mobility changes across space and social groups.}
  
\setlength{\leftskip}{1cm}\textbf{A:} Thanks for this comment, it is very similar to R1.4, therefore, we include here the same response. We agree that multiple covariates may influence both mobility behaviour and recovery dynamics. As such, we analysed variations in these components across deprivation and population density categories. Particularly, we used the Relative Deprivation Index (RDI) because it is a composite index that captures six overlapping dimensions reflecting geographic differences in: (i) built-up density; (ii) child dependency ratio; (iii) infant mortality rate; (iv) human development; (v) nighttime lights intensity; and (vi) nighttime lights slope \parencite{RDI2022}. These dimensions are intended to reflect differences in education, health and infrastructural development \parencite{RDI2022}, mitigating the concern that the analysis omits key structural determinants. 

We also note that (i) our models incorporate a random intercept to account for unobserved structural differences across regions that are not explicitly captured by the included covariates \parencite{gelman2007data}, (ii) our choice of covariates for the models is based on conceptual relevance and on availability of data across countries and at a spatial resolution fine enough to match the granularity of our mobility analysis. 

Within these constraints, we recognise the value of explicitly analysing mobility for areas with varying age structure, as age can plausibly influence mobility patterns \parencite{schwartz1976migration} and capacity to recover, as well as the levels of digital data coverage (as we see in \parencite{cabrera2025bias}). Accordingly, we extended our analysis to control for age groups, as detailed below.

\textcolor{blue}{
\textbf{Action:} We have revised the manuscript taking the following actions:
\begin{itemize}[noitemsep, topsep=0pt]
    \item First, we expanded the modelling framework to incorporate demographic age structure by computing the old-age dependency ratio for each spatial unit, defined as the percentage of residents aged 65 and above relative to the total population. We focus specifically on the old-age dependency ratio, as child dependency is already indirectly captured through the RDI. We derived the old-age dependency ratio from WorldPop data, which provide globally consistent, fine-resolution population estimates to match the spatial resolution of our mobility data. To examine whether demographic composition influences mobility patterns, we categorised spatial units into three groups (low, medium, high) based on their old-age dependency levels. The visual results are collected in the new Supplementary Figure SF11. The statistical models (SF11 and new Supplementary Tables ST4-9) show no significant differences in recovery rates (time slopes) across old-age dependency ratio categories. Intercepts do differ across the three old-age dependency categories, indicating variation in the initial magnitude of the mobility shock, but these differences do not follow a systematic gradient across countries. This contrasts with the reported patterns for population density and RDI categories. For this reason, we focus on population density and RDI for the main analysis. We report the age-structure results in the Supplementary Information.
    \item Second, we have revised the text in the ``Results'' section to (i) describe in greater detail the multidimensional components of the RDI, which capture educational, health and infrastructure characteristics relevant to mobility (further details of the changes in the text can be found in comment R3.7), and (ii) to describe the inclusion of the additional old-age dependency ratio variable in our models.
\end{itemize}
}


\setlength{\leftskip}{0cm} \textit{\textbf{R2.9}: In several figures (e.g., Figures 2 and 8), ``deprived'' and ``dense'' areas are visualized using the same color ramps, implying they are interchangeable or strongly correlated. This conflation is misleading. High population density does not necessarily imply deprivation, and vice versa. I recommend disaggregating these two variables using separate legends or color schemes to improve clarity and avoid interpretive confusion.}

\setlength{\leftskip}{1cm}\textcolor{blue}{
\textbf{Action:} We have revised all the Figures in the manuscript, ensuring that the colour scheme for socioeconomic deprivation and for population density are different.
}

\setlength{\leftskip}{0cm} \textit{\textbf{R2.10}: In Figure 1, the scale does not distinguish between areas with zero population and those with very low population counts. It would be clearer to represent zero-population areas as a separate category.}

\setlength{\leftskip}{1cm}\textbf{A:} Thanks for pointing this out. We think the reviewer was referring to Figures 2, 3 and 4, where the legend includes the minimum and maximum value for each population density category. We would like to clarify that the lowest population category in the analysis of mobility always includes more than 0 persons per sqkm, as areas with 0 people cannot record observed movements (outflows).

\textcolor{blue}{
\textbf{Action:} We have revised the legends in Figures 2, 3 and 4 in the main body of the manuscript and Supplementary Figure SF10. To avoid confusion, we now write $<393$ for the lowest category. In addition, we have added some text in the ``Methods – Classification of tiles by urbanisation and socio-economic deprivation levels'' section to clarify this. Specifically: \textit{``\ul{Population density and deprivation categories are constructed using gridded population estimates and relative deprivation indices. Spatial units with zero population are included in the underlying population layers used to define category thresholds, however, these units are excluded from the mobility analysis, as they cannot generate observed movements. As a result, the lowest population density category used in the mobility analysis includes only spatial units with population density greater than zero.}}''
}

\vspace{10pt}

\setlength{\leftskip}{0cm} \textbf{REVIEWER 3}
\vspace{-16pt}
\begin{center}
    \rule{\textwidth}{.1pt}
\end{center}

\vspace{5pt}

\setlength{\leftskip}{0cm} \textit{\textbf{Reviewer 3}: Thanks for the opportunity to review this paper. The paper focuses on the changes on urban mobility in Latin America after the COVID-19 pandemic.}

\textit{\textbf{R3.1}: Abstract. The abstract describes in general terms the data used for this study, but there is little information regarding data collection and research design. There is also little information regarding modal split. Thus, it is unclear what modal changes the authors are referring to.}

\setlength{\leftskip}{1cm}\textbf{A:} Thanks for the time taken to carefully review our paper. We would like to clarify that our article did not focus on modal split. That is not the aim of the paper. The Meta–Facebook mobility dataset that we used does not contain information on transport modes or modal split. The dataset is derived from mobile phone location records and these do not include modal split information. The Meta–Facebook data capture aggregate movements between spatial tiles inferred from location signals of active Facebook users, rather than trip purpose or transport mode. Thus, our analysis focuses on changes in the overall volume and spatial distribution of mobility across socio-economic and rural–urban gradients.

\textcolor{blue}{
\textbf{Action:} We revised the ``Introduction'' section by explicitly adding the following statement: \textit{\uline{``Although Meta collects individual-level location information, the data made available to researchers consist only of aggregate movement flows between spatial units and do not distinguish transport modes''}}. We also clarify the use of aggregate data in the ``Methods - Data'' subsection within ``Methods'', which is placed after ``Discussion'', as per the journal's formatting requirements. 
}

\setlength{\leftskip}{0cm} \textit{\textbf{R3.2}: Introduction. The introduction section refers to the impacts of the COVID-19 pandemic on urban mobility. However, this is something that requires a more precise description of the phenomena of interest, given that changes occurred due to multiple factors. Government restrictions also implied changes, travel behavior issues as well. I think the manuscript would benefit by including the implications of multiple factors on urban mobility changes due to the pandemic and the policies implemented in the pandemic.}

\setlength{\leftskip}{1cm}\textbf{A:} We agree that multiple interacting factors, including government restrictions, behavioural adaptation and risk perceptions, are all important and potentially contributed to shape mobility behaviour during the COVID-19 pandemic. We discuss this in paragraphs 1 and 3 of our ``Introduction'', and in the newly revised ``Discussion''.  We also incorporate these important considerations in our analysis, as described in the ``Results'' and ``Methods'' section, where we assess the influence of changes in government restrictions on mobility by including the COVID-19 Stringency Index created by \parencite{hale2021global} from the University of Oxford in our regression models. We now clarify this. The index systematically tracks and scores the strictness of containment and closure policies on a 0–100 scale, where higher values indicate more stringent measures. The index aggregates daily data on nine policy indicators including school and workplace closures, cancellation of public events, restrictions on gatherings and movement, public transport closures, stay-at-home requirements, international travel controls and public information campaigns. The detail of the various analyses we conducted using the Stringency Index as a covariate are reported in ST5, ST7, ST9, ST11, ST13, ST15, ST17, ST19 and ST21.

\textcolor{blue}{
\textbf{Action:} We have taken two actions. 
\begin{itemize}[noitemsep, topsep=0pt]
    \item First, we have revised the text in the ``Results'' section to expand the description of the influence of restrictive measures on mobility as suggested by the reviewer, and how we incorporate this into the analysis. In particular, we have added the following text towards the end of the first paragraph of the section ``Results - Early mobility declines during the pandemic shape long-term levels of mobility'': \textit{``We modelled the trend component as a function of time, stringency levels and attributes of the places of origin, including population density, relative deprivation and old-age dependency ratio (the latter category reported only in the Supplementary Information), in a mixed-effects modelling framework – see Methods section and ST3 for details. \ul{Stringency levels are measured via the COVID-19 Stringency Index created by} \parencite{hale2021global}\ul{, which we include to assess the influence of changes in government restrictions on mobility. The index aggregates daily data on nine policy indicators including school and workplace closures, cancellation of public events, restrictions on gatherings and movement, public transport closures, stay-at-home requirements, international travel controls and public information campaigns.}''}
    \item Second, we have revised the ``Discussion'' carefully, to better highlight the policy and behavioural implications of the changes in mobility that we evidence through our analysis. In particular, we have added the following two paragraphs: \textit{\ul{``The evidence reported here potentially has important social equity implications by exacerbating pre-existing disparities in Latin America. The region is among the most unequal regions in the world} \parencite{WIReport22}\ul{. The bottom 50\% of the Latin American population holds only around 10\% of wealth} \parencite{WIReport22}\ul{, and is likely to face a disproportionately high social burden due to a more limited capacity to work remotely and a continued reliance on physical mobility for accessing employment and services. This entails financial and non-financial opportunity costs, such as transport fares and travel time. By contrast, the top 10\% of the population, which holds approximately 55\% of wealth} \parencite{WIReport22}\ul{, may continue to reduce travel due to greater flexibility in work arrangements and access to digital alternatives. As a result, uneven mobility recovery across socio-economic groups may function as a mechanism through which pandemic-era disruptions become embedded into long-standing patterns of inequality and spatial mismatches between residence and opportunity, with potential consequences for labour market participation, educational attainment and access to essential services. These sustained mobility changes also have important implications for Latin American cities and their governance. Sustained suppression of mobility into urban cores risks undermining commercial activity and overall urban vibrancy, particularly if higher-income populations continue to travel less frequently to central areas. }}
\end{itemize} }

\setlength{\leftskip}{1.3cm}\textcolor{blue}{\textit{\ul{These dynamics may reinforce already emerging patterns of urban centre decline in some large Latin American cities, as we saw in European cities during the 1970s and 1980s} \parencite{THORSTEN}\ul{. In Latin America, where metropolitan governance is often fragmented across municipalities and fiscal capacity for large-scale interventions remains limited} \parencite{un2022world, inostroza2013urban, fay2017rethinking, irazabal20governance}\ul{, the persistence of depressed urban core activity highlights the need for targeted, context-specific responses, rather than broad calls for inclusive or sustainable mobility. Policies prioritising affordable and reliable public transport connections to central employment and service hubs, particularly during peak hours and along high-demand corridors, may be more effective than uniform system-wide measures} \parencite{rivas2019urban, guzman2022effects}\ul{. More broadly, the observed heterogeneity in recovery trajectories underpins the importance of metropolitan-scale coordination in transport and urban planning, as city-level responses alone may be insufficient to address persistent shifts in urban activity patterns} \parencite{duque2021institutional}\ul{. In this scenario, near real-time mobility data offer a valuable tool for identifying emerging signals of urban decline and supporting adaptive, evidence-based policy interventions in contexts where traditional data sources remain delayed or incomplete. ''}}.}

\setlength{\leftskip}{0cm} \textit{\textbf{R3.3}: Literature review. The document lacks a literature review section. There are already several studies looking at urban mobility during the pandemic of COVID-19. I think the document would benefit from the inclusion of a literature review section. If the editors of the journal prefer that the literature review should be part of the introduction section, the manuscript would benefit from the inclusion of the literature review in the introduction section.}

\setlength{\leftskip}{1cm}\textbf{A:} We thank the reviewer for this thoughtful comment. We have expanded the ``Introduction'' to include relevant studies assessing the impact of the COVID-19 pandemic on population mobility globally, with particular attention to Latin American countries. In line with the scope and style of \textit{Nature Communications}, we focus on synthesising the key findings from the existing literature rather than providing an exhaustive review, in order to keep the manuscript concise (maximum 5,000 words) and to emphasise the novel contributions of our study.

We also note that the journal’s formatting guidelines specify that \textit{``Articles should prioritise a focused Introduction that provides the necessary background for the study, followed by Results, Discussion and Methods sections.''} Accordingly, the inclusion of a full literature review section is not standard for this journal. We believe that the revised ``Introduction'' now provides sufficient context while adhering to the journal’s editorial requirements.

\textcolor{blue}{
\textbf{Action:} We have revised the ``Introduction'' to add a new paragraph that includes a more extensive list of key references relevant to our work. Particularly, we have added: \textit{``Data resulting from social interactions on digital platforms have emerged as a unique source of information to capture human population movement in less developed countries \parencite{rowe2023big}. Particularly, location data from mobile phone applications have become a prominent source to sense patterns of human mobility at higher geographical and temporal resolution in real time \parencite{calafiore2023, iradukunda2025}\ul{. These data sources have consistently documented the sharp and widespread mobility declines during the early months of the pandemic across Latin American cities} \parencite{RoweCCA24, covidMexico24, elejalde2024, demography25}\ul{, providing a well-established empirical foundation for focusing on the persistence and inequality of the post-pandemic recovery in the present study.''}}
}

\setlength{\leftskip}{0cm} \textit{\textbf{R3.4}: Research design and methods. The paper includes research questions in the introduction section. However, there is not a clear description of the research design and methods.}

\setlength{\leftskip}{1cm}\textbf{A:} We would like to note that the paper provides an extensive description of the research design and methods in the ``Methods'' section, which is included after the ``Discussion'' section, as per the journal formatting requirements. The ``Methods'' section offers all the details of the analysis, including alternative model specifications, techniques and references to the various sections of Supplementary Information, where we have included extensive validation and robustness checks. 

\textcolor{blue}{
\textbf{Action:} We have revised the text to more clearly indicate readers where they can find details of the research design and methods used in the paper and in the Supplementary Information. Particularly, we have included additional text at the beginning of the ``Results'' section, specifying the following: \textit{``\ul{Details about our approach can be found in the Methods section}''}. And also, in the caption of the newly revised Figure 1, which offers an overview of the methodology: \textit{\ul{``Overview of our methodology. Detailed descriptions are provided in the Methods section. [...]}''}. Further details on our updates to Figure 1 and its caption can be found in comment R1.6 and R3.5.
}

\setlength{\leftskip}{0cm} \textit{\textbf{R3.5}: Results. The paper includes in figure 1 the description of data processing and analysis. I think authors should cite the references that support the research methods. There is no description of travel modes in the data analysis. }

\setlength{\leftskip}{1cm}\textbf{A:} Thanks for the opportunity to clarify. Figure 1 provides an overview of the methodology. Full details of the methods can be found in the ``Methods'' section, where we have included all the relevant references. The ``Methods'' section is placed after the ``Discussion'', following the journal's formatting requirements. We have revised Figure 1, following also the advice from another reviewer, simplifying its content for a clearer overview of our methodology. We have included a more detailed caption for this Figure, which refers readers to the ``Methods'' section for further details.

The paper does not aim to analyse modal split. It seeks to assess the long-term influence of the COVID-19 pandemic on population movements. We have clarified this in the ``Introduction''.

\textcolor{blue}{
\textbf{Action:} We have taken three actions.
\begin{itemize}[noitemsep, topsep=0pt]
    \item First, we revised Figure 1, following also the advice from another reviewer (see comment R1.6). The figure is now simpler and reflects better the overview of our methodology. Importantly, we have updated the caption, to refer the readers to the full description of the methodology in the ``Methods'' section, where references can also be found. Particularly, the caption now reads as: \textit{\ul{``Overview of our methodology. Detailed descriptions are provided in the Methods section. Facebook mobility and population data (Meta), relative deprivation (NASA), population estimates (WorldPop) and functional urban area boundaries (Global Human Settlement Layer) are used as inputs. Mobility data are calibrated to reduce spatial and temporal biases and analysed within each country using statistical models to estimate urban mobility recovery. Validation assesses consistency with WorldPop populations and robustness across calibration and model specifications.''}}.
    \item We have added text in the introduction to clarify that we do not analyse model split in the paper, as this information is not available in our data. Particularly, as mentioned in our response to R3.1, we added \textit{\uline{``Although Meta collects individual-level location information, the data made available to researchers consist only of aggregate movement flows between spatial units and do not distinguish transport modes''}}.
    \item We have added a paragraph in the ``Discussion'', elaborating on the limitations of the Meta-Facebook data, among which we highlight the lack on information on model split. See comment R3.11 for further details on the changes we made to the ``Discussion''.
\end{itemize}
}

\setlength{\leftskip}{0cm} \textit{\textbf{R3.6}: The presentation in figure 2 of the data from a descriptive approach is based on several assumptions that are not clearly explained. For instance, population density differs depending on the definition of what is urban on each country, thus, it is unclear how the analysis incorporates with precision urban population density.}

\setlength{\leftskip}{1cm}\textbf{A:} Thank you for the comment. We describe all the key assumptions of the methodology in the ``Methods'' section, which is placed after ``Discussion'', following the journal's formatting requirements. This includes the classification of areas by population density and relative deprivation. There is a dedicated section for this called ``Methods - Classification of tiles by urbanisation and socio-economic deprivation levels''. In particular, the population density classification alluded in the comment above are built from a harmonised, gridded population surface of 1km obtained from WorldPop. We aggregated individual grids to the Bing tiles used by Meta and classified the resulting tiles into five density classes using fixed, cross-country thresholds. These classes represent areas along the rural-urban continuum: the highest density areas characterise urban centres, such as Central Business Districts (CBDs); followed by dense urban areas; semi-dense and suburban areas; rural areas; and, low-density rural areas. Broadly speaking, these areas correspond to the areas identified in the degree of urbanisation.

\textcolor{blue}{
\textbf{Action:} We have revised the text and relevant sections of the paper to indicate readers where they can find details of the research design and methods used in the paper and Supplementary Information. Particularly, at the beggining of the ``Results'' section, we have added: \textit{``We analysed post-COVID-19 mobility patterns over time relative to pre-pandemic baseline levels, using anonymised location data from Meta-Facebook users in Argentina, Chile and Colombia, combined with high-resolution geospatial datasets (see Figure 1). \ul{Details about our approach can be found in the Methods section}''}. And at the end of the first paragraph of ``Results'' we have also added \textit{\uline{``Further details of this classification can be found in Methods, `Classification of tiles by urbanisation and socio-economic deprivation levels'.''}}
}

\setlength{\leftskip}{0cm} \textit{\textbf{R3.7}: The calculation of the deprivation index is not clear. I think the paper should include a table with all variables used in the analysis, and then, it is important to show the descriptive statistics of those variables.}

\setlength{\leftskip}{1cm}\textbf{A:} We'd like to thank the reviewer for the opportunity to clarify. We used the Relative Deprivation Index (RDI) created by NASA and described here: \parencite{RDI2022}. The Relative Deprivation Index (RDI) is a composite index that captures six overlapping dimensions reflecting geographic differences in: (i) built-up density; (ii) child dependency ratio; (iii) infant mortality rate; (iv) human development; (v) nighttime lights intensity; and, (vi) nighttime lights slope \parencite{RDI2022}. These dimensions are intended to reflect differences in education, health and infrastructural development \parencite{RDI2022}.

A detailed description of all the data sources used in the analysis can be found in the ``Methods - Data'' section. We have also included the data sources in the newly revised Figure 1 and its updated caption.

\textcolor{blue}{
\textbf{Action:} We have revised the text to provide more details about the RDI and included relevant descriptions of the covariates used in our analysis. Particularly we have take the following actions:
\begin{itemize}[noitemsep, topsep=0pt]
    \item First, we have revised the text to provide more details about the RDI in the ``Results'' section. Particularly, we have added the following: \textit{``We defined [...] three socioeconomic deprivation deciles \ul{based on NASA's Relative Deprivation Index (RDI). The RDI captures multiple structural dimensions, including built-up density, child dependency, infant mortality, subnational human development, and night-time lights intensity and slope, reflecting educational, health and infrastructural characteristics relevant to mobility. [...]. The resulting deprivation categories represent low, medium and high deprivation areas}''}.
    \item Second, we have redesigned SF7-9 in the Supplementary Information to include more descriptive statistics of the covariates included in the model. Particularly, SF7 now provides descriptive statistics of two key covariates in the analysis, i.e. population density (for urban-rural classification) and relative deprivation index (for socio-economic classification). SF8 provides descriptive statistics of the population represented by the Meta-Facebook mobility data. SF9 shows the joint distribution of spatial units by population density and relative deprivation index categories, providing descriptive context for the combined socioeconomic and density structure of the study areas. We also included notes in the Figure captions stating how they link to the results reported in the main manuscript. Further details of all the data sources used in the analysis can be found in the ``Methods - Data'' section.
    \item Third, we have updated Figure 1, to more clearly reflect all the input data sources in our analysis. Details on how we updated this Figure and its caption can be found in comment R1.6 and R3.5.
\end{itemize}
}

\setlength{\leftskip}{0cm} \textit{\textbf{R3.8}: The random effects models are shown by country, but there is not any analysis looking at changes over time.}

\setlength{\leftskip}{1cm}\textbf{A:} A key aim of our paper is to analyse the temporal persistence of changes in mobility due to COVID-19. All our figures, including our mixed-effects model estimates, report evidence on the recorded changes in mobility over time. As described in the paper, mobility trends are modelled as a function of time based on both linear and nonlinear components, while accounting for random intercepts and slopes at subnational levels. The coefficients derived for each country thus represent the temporal evolution of mobility and its interaction with socio-economic and population density gradients.

\textcolor{blue}{
\textbf{Action:} We have clarified the text to describe the ways in which we modelled mobility patterns as a function of time. This is explained in various points in the paper. Particularly, in the ``Results - Early mobility declines during the pandemic shape long-term levels of mobility'' section, we have revised the following paragraph for increased clarity: \textit{``Discrepancies in the resilience of areas may underpin the observed differences in mobility levels. To quantify the resilience of areas, \ul{we model the trends of mobility recovery to pre-pandemic levels as a function of time}. To this end, we decomposed the time series of the percentage change in mobility into trend, seasonal and residual components using a moving-average–based method. Robustness checks using alternative trend extraction techniques are reported in Supplementary Table (ST) ST2. The time series are displayed in SF10. \ul{We modelled the trend component as a function of time, stringency levels and attributes of the places of origin, including population density, relative deprivation and old-age dependency ratio} (the latter category reported only in the Supplementary Information), in a mixed-effects modelling framework – see Methods section and ST3 for details. [...] Figure 3 (a) and (b) reports estimates from \ul{time-only models}, which were also favoured by the Akaike Information Criterion (AIC). [...]''}. Further details on the methods used for modelling mobility changes as a function of time are available in ``Methods - Trend analysis'' and full results can be found revised tables ST4-21 from Supplementary Information.
}

\setlength{\leftskip}{0cm} \textit{\textbf{R3.9}: The analysis shown in figure 4 is interesting, but it is not clear how population density was estimated, while there is not any information regarding travel modes.}

\setlength{\leftskip}{1cm}\textbf{A:} Thanks for the opportunity to clarify. We describe all the methodological decisions taken in the ``Methods'' section. This includes the classification of areas by population density and their related modelling estimates. In particular, the population density classification alluded in the comment above are built from a harmonised, gridded population surface of 1km obtained from WorldPop. We aggregated individual grids to the Bing tiles used by Meta and classified the resulting tiles into five density classes using fixed, cross-country thresholds. These classes represent areas along the rural-urban continuum: the highest density areas characterise urban centres, such as Central Business Districts (CBDs); followed by dense urban areas; semi-dense and suburban areas; rural areas; and low-density rural areas. We used these categories as covariates in our mixed-linear regression models to derive population density estimates for each category.

Additionally, we would like to highlight that our article did not focus on modal split. That is not the aim of the paper. The Meta–Facebook mobility dataset that we used does not contain information on transport modes or modal split. The dataset is derived from mobile phone location records and these do not include modal split information. The Meta–Facebook data capture aggregate movements between spatial tiles inferred from location signals of active Facebook users, rather than trip purpose or transport mode. Thus, our analysis focuses on changes in the overall volume and spatial distribution of mobility across socio-economic and rural–urban gradients.


\textcolor{blue}{
\textbf{Action:} Thanks. We revised the ``Introduction'' section to explicitly state that the Facebook data represent aggregate human movement between spatial units and do not identify transport modes. Particularly, as described in comment R3.1, we added the following in the ``Introduction'': \textit{\uline{``Although Meta collects individual-level location information, the data made available to researchers consist only of aggregate movement flows between spatial units and do not distinguish transport modes''}}. We also clarify the use of aggregate data in the ``Methods - Data'' subsection. Furthermore, we have revised the entire text of the paper to clearly indicate readers where they can find details of the research design and methods used in the paper and Supplementary Information.
}

\setlength{\leftskip}{0cm} \textit{\textbf{R3.10}: Methods. The methods section is presented after the results section. I think the order should be the opposite. There is little information regarding how Facebook data has limitations. Facebook data is not an origin destination survey, and I wonder if authors could describe travel modes as well as trip purpose based on the data. The authors do not include information regarding the population characteristics, which might be useful in the analysis. The description of research methods does not mention how there is analysis between and within countries, and how the research designs work spatial issues.}

\setlength{\leftskip}{1cm}\textbf{A:} We would like to note that the paper is structured according to the journal's guidelines; that is, the ``Results'' Section being placed before the ``Methods'' section.  We would be happy to move the ``Methods'' before ``Results'' if the Editors find it appropriate during the production stage if the paper reaches that stage.

As described above, our article does not focus on modal split or trip purpose, as this information is not available in the Facebook mobility dataset. Facebook mobility data are anonymised and aggregated, and therefore do not constitute a traditional origin–destination survey. Data on travel mode, trip purpose, or individual-level socioeconomic characteristics is not included. We have clarified these limitations in the revised manuscript.

Regarding population characteristics, while individual-level attributes are not available due to data anonymisation, we have included a new Supplementary Figure (SF8) in the Supplementary Information, which shows the distribution of Facebook users across population groups based on the urban–rural and the relative deprivation classification of areas. This addition helps contextualise the representativeness of the dataset.

Finally, each country is analysed independently, with the same methodological approach applied across all cases to identify internal spatial patterns, rather than to compare countries directly. We have clarified this in the Introduction.

\textcolor{blue}{
\textbf{Action:} We have taken three actions. \begin{itemize}[noitemsep, topsep=0pt]
    \item First, we have added text in the ``Introduction'' to note that the focus is not on the modal split. See comment R3.1 and R3.5 for details of the changes we incorporated to clarify this. We have also added a paragraph in the ``Discussion'' noting the limitations in our data. See comment R3.11 for details.
    \item We have added Supplementary Figure 8 to give a better idea of the population groups represented by our data. We also refer to this Figure in the main body of the text, particularly in the ``Methods'' section.
    \item We now clarify in the ``Introduction'' that our analysis focuses on within-country movements only. Particularly, we added the following text \textit{``Drawing on a large dataset of 170 million observations from Meta-Facebook users' mobile location data, we aim to assess the extent and pace of urban mobility recovery \uline{within countries} in Argentina, Chile and Colombia over a 26-month period from April 2020 to May 2022 following abrupt declines during early phases of the COVID-19 pandemic. \uline{Each country is analysed independently, focusing on internal mobility patterns.}''}.
\end{itemize}}

\setlength{\leftskip}{0cm} \textit{\textbf{R3.11}: Discussion. The paper lacks a clear description of the limitations of the research design and analysis.}

\setlength{\leftskip}{1cm}\textbf{A:} This is a fair point. We have added text to explicitly point to some key limitations of our analysis.

\textcolor{blue}{
\textbf{Action:} We have revised the ``Discussion'' to include more explicit information regarding the challenges in our methodology and data. Particularly, we have added the following paragraph: \textit{\ul{``Despite our robustness and validation tests, key challenges should be considered when interpreting our results. First, limitations related to the use of aggregated mobile-phone–derived mobility indicators. These data are generated from a non-random sample of users providing variable coverage across space and time, and potentially under-representing population groups with lower smartphone access, limited digital connectivity or higher privacy opt-out rates} \parencite{cabrera2025bias}\ul{. As such, inferred mobility changes may not fully reflect the behaviour of all sociodemographic groups, particularly in more disadvantaged areas. Second, the pre-pandemic baseline used to define reference mobility levels is limited to a 45-day period provided by Meta and therefore does not capture longer-term seasonal variability. While this baseline precedes mobility restrictions and accounts for day-of-week effects, it may affect estimates of the absolute magnitude of the initial mobility drop, though relative recovery trajectories are less sensitive to this limitation. Third, our data do not provide information on modal split, limiting inference on the persistence of changes in specific transport modes, such as walking or cycling} \parencite{zhang2023impacts}\ul{. Fourth, observed mobility changes reflect outcomes rather than travel motivations, constraining our ability to distinguish behavioural adaptation from mobility reductions driven by affordability, service availability or labour informality. Finally, our analyses are observational and conducted at an aggregate spatial scale, so associations between mobility recovery and city characteristics should not be interpreted as causal and may be influenced by unobserved factors, such as changes in transport supply or concurrent economic shocks. These limitations motivate future work combining digital trace mobility data with external validation sources and complementary data to assess representativeness, modal shifts and causal mechanisms.''}}
}

\newpage

\clearpage
\pagestyle{plain}   
\printbibliography  


\end{document}

